% ============================================================
%  DISCOVERY CURRICULUM — Computation Journey 1, Module 1
%  "The Nature of Information"
%  Compile with: pdflatex module 1.tex  (run twice for references)
% ============================================================\documentclass[12pt, a4paper]{article}

\usepackage[a4paper, top=2.8cm, bottom=2.8cm, left=3.0cm, right=2.5cm, includeheadfoot, headheight=15pt]{geometry}
\usepackage{amsmath, amsthm}
\usepackage[charter, cal=cmcal]{mathdesign}
\usepackage[scaled=0.92]{helvet}
\usepackage{xcolor}
\usepackage{tcolorbox}
\tcbuselibrary{skins, breakable, theorems}
\usepackage{enumitem}
\usepackage{graphicx}
\usepackage{tikz}
\usetikzlibrary{arrows.meta, positioning, calc, trees}
\usepackage{booktabs}
\usepackage{array}
\usepackage{fancyhdr}
\usepackage{titlesec}
\usepackage{parskip}
\usepackage{microtype}
\usepackage{hyperref}
\hypersetup{colorlinks=true, linkcolor=CompNavy, urlcolor=CompBlue}

% ── Colours (Sleek Slate/Indigo Palette) ────────────────────
\definecolor{CompNavy}{HTML}{0F172A}    % Slate 900
\definecolor{CompBlue}{HTML}{2563EB}    % Blue 600
\definecolor{CompGold}{HTML}{D97706}    % Amber 600
\definecolor{CompLightBlue}{HTML}{F0F6FF}
\definecolor{CompLightGold}{HTML}{FEF3C7}
\definecolor{DarkGrey}{HTML}{334155}    % Slate 700
\definecolor{BoxGreen}{HTML}{F0FDF4}    % Green 50
\definecolor{SidebarGrey}{HTML}{F8FAFC} % Slate 50

% ── Section formatting ───────────────────────────────────────
\titleformat{\section}{\Large\bfseries\color{CompNavy}}{}{0em}{}[{\color{CompBlue}\vspace{2pt}\rule{\linewidth}{1.5pt}}]
\titleformat{\subsection}{\large\bfseries\color{CompBlue}}{}{0em}{}
\titleformat{\subsubsection}{\normalsize\bfseries\color{CompGold}}{}{0em}{}

% ── Header / Footer ─────────────────────────────────────────
\pagestyle{fancy}
\fancyhf{}
\fancyhead[L]{\small\color{DarkGrey}\textit{Discovery Curriculum}}
\fancyhead[R]{\small\color{DarkGrey}\textit{Computation Journey 1 $\cdot$ Module 1}}
\fancyfoot[C]{\small\color{CompBlue}\thepage}
\renewcommand{\headrulewidth}{0.4pt}

% ── Environments (tcolorbox Unification) ─────────────────────

% Puzzle box — opening problem
\newtcolorbox{puzzlebox}[1][]{
  enhanced, breakable,
  colback=CompNavy!5, colframe=CompNavy,
  boxrule=0pt, leftrule=2.5pt,
  sharp corners,
  left=12pt, right=10pt, top=10pt, bottom=10pt,
  before skip=14pt, after skip=10pt,
  #1
}

% Discovery box — moment of crystallisation
\newtcolorbox{discoverybox}[1][]{
  enhanced, breakable,
  colback=CompLightGold, colframe=CompGold,
  boxrule=0pt, leftrule=2.5pt,
  sharp corners,
  left=12pt, right=10pt, top=10pt, bottom=10pt,
  before skip=14pt, after skip=10pt,
  #1
}

% Name box — formal definition reveal
\newtcolorbox{namebox}[1][]{
  enhanced, breakable,
  colback=CompLightBlue, colframe=CompBlue,
  boxrule=0.8pt,
  arc=4pt,
  left=12pt, right=10pt, top=10pt, bottom=10pt,
  before skip=14pt, after skip=10pt,
  #1
}

% Reflection box
\newtcolorbox{reflectionbox}[1][]{
  enhanced, breakable,
  colback=BoxGreen, colframe=DarkGrey!35,
  boxrule=0.8pt,
  arc=4pt,
  left=12pt, right=10pt, top=10pt, bottom=10pt,
  before skip=14pt, after skip=10pt,
  #1
}

% Thought provocation — inline philosophical aside
\newtcolorbox{thought}[1][]{
  enhanced, breakable,
  colback=SidebarGrey, colframe=CompGold,
  boxrule=0pt, leftrule=1.5pt,
  sharp corners,
  left=10pt, right=10pt, top=7pt, bottom=7pt,
  before skip=10pt, after skip=10pt,
  #1
}

% Problem counter
\newcounter{problem}
\newenvironment{problem}[1][]{%
  \refstepcounter{problem}%
  \par\vspace{6pt}%
  \noindent\textbf{\color{CompBlue}Problem \theproblem}%
  \ifx&#1&\else\ \textit{(#1)}\fi\textbf{.}\enskip
}{\par\vspace{4pt}}

% Deep problem counter
\newcounter{deepproblem}
\newenvironment{deepproblem}[1][]{%
  \refstepcounter{deepproblem}%
  \par\vspace{6pt}%
  \noindent\textbf{\color{CompNavy}$\bigstar$\ Deep Problem \thedeepproblem}%
  \ifx&#1&\else\ \textit{(#1)}\fi\textbf{.}\enskip
}{\par\vspace{4pt}}

% Exercise counter
\newcounter{exercise}
\newenvironment{exercise}[1][]{%
  \refstepcounter{exercise}%
  \par\vspace{6pt}%
  \noindent\textbf{\color{CompGold}Exercise \theexercise}%
  \ifx&#1&\else\ \textit{(#1)}\fi\textbf{.}\enskip
}{\par\vspace{4pt}}

% ============================================================
\begin{document}
% ============================================================

\thispagestyle{empty}
\begin{center}
{\small\color{CompBlue}\textbf{DISCOVERY CURRICULUM}\quad$\cdot$\quad\textbf{JOURNEY 1: NATURE OF INFORMATION}}

\vspace{0.6cm}
{\color{CompNavy}\rule{\linewidth}{2pt}}
\vspace{0.5cm}

{\Huge\bfseries\color{CompNavy}Module 1.1}\\[0.3cm]
{\LARGE\bfseries\color{CompBlue}The Telegraph Puzzle: Signals and Representation}\\[0.4cm]
{\color{CompGold}\rule{0.4\linewidth}{1.2pt}}\\[0.4cm]
{\large\itshape\color{DarkGrey}
Before we compute, we must represent.\\
How do we measure the raw material of thought?}

\vspace{1cm}
\begin{tabular}{rl}
\color{DarkGrey}\textbf{Journey}           & Journey 1 --- The Nature of Information \\[4pt]
\color{DarkGrey}\textbf{Module}            & 1 of 5 \\[4pt]
\color{DarkGrey}\textbf{Prerequisites}     & Combinatorics, basic probability, real analysis limits \\[4pt]
\color{DarkGrey}\textbf{You will discover} & Prefix-free codes, Kraft's Inequality, Huffman coding, Entropy \\[4pt]
\color{DarkGrey}\textbf{Format}            & 8 problems $+$ 3 solved exercises $+$ 2 deep problems \\
\end{tabular}

\vspace{1.5cm}
{\color{CompNavy}\rule{\linewidth}{2pt}}
\end{center}

\newpage


% ============================================================
\section{Part 1 \texorpdfstring{---}{--} The Puzzle}
% ============================================================

\begin{puzzlebox}
\textbf{The Telegraph Puzzle.}

You are a telegraph operator in 1844. You can transmit only two types of electrical pulses along a copper wire: a short pulse (\textbf{0}) and a long pulse (\textbf{1}). 

You want to send messages consisting of letters from a simplified 4-letter alphabet: $\Sigma = \{A, B, C, D\}$. 

\begin{enumerate}[label=\textbf{\Alph*.}]
  \item If you represent each letter with a fixed-length code ($A \mapsto 00$, $B \mapsto 01$, $C \mapsto 10$, $D \mapsto 11$), sending a sequence of letters is straightforward. The stream $00011110$ is uniquely parsed as $A-B-D-C$.
  \item However, you observe that letter $A$ appears 70\% of the time, $B$ appears 15\%, $C$ appears 10\%, and $D$ appears only 5\%. To save time and copper wire, you want to make common letters shorter.
  \item You propose the variable-length code:
  \[
    A \;\mapsto\; 0, \qquad B \;\mapsto\; 1, \qquad C \;\mapsto\; 01, \qquad D \;\mapsto\; 10.
  \]
\end{enumerate}

\textbf{The Breakdown:} A message arrives: $010$. The receiver is confused. Is this $A-D$ ($0$ followed by $10$)? Is it $C-A$ ($01$ followed by $0$)? Or is it $A-B-A$ ($0$ followed by $1$ followed by $0$)? 

\medskip
\textbf{Core Questions:}
\begin{itemize}
  \item Why did the fixed-length code parse uniquely, while the shorter variable-length code failed?
  \item How can a variable-length code be structured to be uniquely decodable without delimiters?
  \item What is the mathematical lower bound on the average number of binary signals required to represent a source?
\end{itemize}
\end{puzzlebox}

\bigskip

% ============================================================
\section{Part 2 \texorpdfstring{---}{--} Exploration}
% ============================================================

\subsection{Step 1: The Geometry of Unambiguous Codes}

To prevent confusion, we need a code where no codeword is a "prefix" (the beginning) of another codeword. Such codes are called \textbf{prefix-free codes}.

\begin{problem}[Prefix-Free Codes as Binary Trees]
Consider a binary tree where branching left represents writing a \textbf{0} and branching right represents writing a \textbf{1}.
\begin{enumerate}[label=(\alph*), topsep=2pt]
  \item Map the confused code from the puzzle ($A \mapsto 0$, $B \mapsto 1$, $C \mapsto 01$, $D \mapsto 10$) onto a binary tree. Where do the letters lie? (Are some descendants of others?)
  \item If we restrict our letters to lie \emph{only} on the leaves (end-nodes) of the tree, explain why no codeword can be a prefix of another.
  \item Construct a leaf-only binary tree for our alphabet $\{A, B, C, D\}$. Write down the corresponding prefix-free codewords.
\end{enumerate}
\end{problem}

\begin{center}
\begin{tikzpicture}[level distance=1.5cm,
  level 1/.style={sibling distance=3cm},
  level 2/.style={sibling distance=1.5cm}]
  \node[circle,draw] {}
    child {node[circle,draw] {}
      child {node[circle,draw,fill=CompLightBlue] {A} edge label={node[above left] {0}}}
      child {node[circle,draw,fill=CompLightBlue] {B} edge label={node[above right] {1}}}
      edge label={node[above left] {0}}
    }
    child {node[circle,draw] {}
      child {node[circle,draw,fill=CompLightBlue] {C} edge label={node[above left] {0}}}
      child {node[circle,draw,fill=CompLightBlue] {D} edge label={node[above right] {1}}}
      edge label={node[above right] {1}}
    };
\end{tikzpicture}
\end{center}

\subsection{Step 2: The Kraft Inequality}

Each choice of codeword length limits the availability of future codewords. If we choose a very short codeword (like $0$), we eliminate the entire left half of the tree. This structural constraint is captured by a beautiful inequality.

\begin{problem}[Discovering the Constraint]
Suppose a prefix-free binary code has codeword lengths $\ell_1, \ell_2, \ldots, \ell_n$.
\begin{enumerate}[label=(\alph*), topsep=2pt]
  \item Let the maximum length of any codeword be $L$. Imagine a full binary tree of depth $L$, which has $2^L$ leaves. If a codeword has length $\ell_i$, how many leaves at depth $L$ are descendants of this codeword (and thus blocked from being used)?
  \item Since the total number of leaves at depth $L$ is $2^L$, write down an inequality summing the blocked leaves for all $n$ codewords.
  \item Divide both sides of your inequality by $2^L$. What simple, depth-independent constraint on the lengths $\ell_i$ emerges?
\end{enumerate}
\end{problem}

\begin{discoverybox}
\textbf{The Kraft Inequality.}
For any prefix-free code over a binary alphabet with codeword lengths $\ell_1, \ell_2, \ldots, \ell_n$, we must have:
\[
  \sum_{i=1}^{n} 2^{-\ell_i} \;\le\; 1.
\]
Conversely, if a set of lengths satisfies this inequality, a prefix-free code with these lengths is guaranteed to exist.
\end{discoverybox}

\subsection{Step 3: Finding the Optimal Code}

Now we want to find a code that minimizes the average length of a message:
\[
  L_{\text{avg}} \;=\; \sum_{i=1}^{n} p_i \ell_i.
\]

\begin{problem}[The Huffman Algorithm]
Let us discover a greedy strategy to build the optimal tree.
\begin{enumerate}[label=(\alph*), topsep=2pt]
  \item If you have two symbols with the lowest probabilities, where should they be placed in the binary tree relative to each other to minimize $L_{\text{avg}}$? (Should they be siblings at the maximum depth?)
  \item Suppose we combine the two least probable symbols into a single "meta-symbol" whose probability is the sum of the two. If we find the optimal code for this smaller set, how do we get the code for the original set?
  \item Use this method to find the optimal prefix-free code for the probabilities:
  \[
    p(A) = 0.70, \qquad p(B) = 0.15, \qquad p(C) = 0.10, \qquad p(D) = 0.05.
  \]
  Calculate its average length $L_{\text{avg}}$. Compare it to the fixed-length average of 2.0.
\end{enumerate}
\end{problem}

\subsection{Step 4: The Fundamental Limit of Information}

Can we compress the message indefinitely, or is there a hard limit? Let us use real analysis to find a lower bound on $L_{\text{avg}}$ subject to Kraft's constraint.

\begin{problem}[Optimizing with Lagrange Multipliers]
Let us treat the lengths $\ell_i$ as continuous real variables. We want to minimize $L_{\text{avg}} = \sum p_i \ell_i$ subject to the equality constraint $\sum 2^{-\ell_i} = 1$.
\begin{enumerate}[label=(\alph*), topsep=2pt]
  \item Write the constraint as $g(\ell_1, \ldots, \ell_n) = \sum e^{-\ell_i \ln 2} - 1 = 0$. Using Lagrange multipliers, set up the Lagrangian:
  \[
    \mathcal{L}(\ell_1, \ldots, \ell_n, \lambda) \;=\; \sum_{i=1}^n p_i \ell_i \;+\; \lambda \left(\sum_{i=1}^n 2^{-\ell_i} - 1\right).
  \]

  \item Differentiate $\mathcal{L}$ with respect to $\ell_i$ and set it to zero. Show that the optimal lengths satisfy:
  \[
    \ell_i^* \;=\; -\log_2 p_i.
  \]
  \item Substitute these optimal lengths back into the formula for $L_{\text{avg}}$. Write down the resulting expression. This quantity is the absolute limit of information density.
\end{enumerate}
\end{problem}

\begin{thought}
\textbf{The fractional bit paradox.}
The optimal length $\ell_i^* = -\log_2 p_i$ is rarely an integer. For instance, if $p_i = 0.7$, then $\ell_i^* \approx 0.51$ bits. We cannot send half a pulse. This means that for a single letter, our actual average length will be slightly higher than the theoretical limit. But if we group letters into blocks (e.g., encoding pairs of letters like $AA, AB, \ldots$), the average length per letter approaches the limit.
\end{thought}

% ============================================================
\section{Part 3 \texorpdfstring{---}{--} Discovery}
% ============================================================

\begin{discoverybox}
\textbf{What have we built?}
1. **Information has a physical limit:** We cannot represent a source of data with fewer bits on average than its *entropy*, no matter how clever our coding scheme.
2. **Prefix-free structures are trees:** The Kraft inequality reveals that coding is a resource allocation problem on a tree: choosing a short codeword blocks a large branch of possibilities.
3. **The Optimal Code is Constructive:** The Huffman algorithm provides a simple, recursive way to find the absolute best prefix code for any known probability distribution.
\end{discoverybox}

% ============================================================
\section{Part 4 \texorpdfstring{---}{--} Naming the Idea}
% ============================================================

\begin{namebox}
\textbf{What mathematicians call it: Shannon Entropy.}

In 1948, **Claude Shannon** (1916–2001) published *"A Mathematical Theory of Communication"*, founding the field of Information Theory. He defined the **entropy** $H(X)$ of a discrete random variable $X$ with probability distribution $P = (p_1, p_2, \ldots, p_n)$ as:
\[
  H(X) \;=\; -\sum_{i=1}^{n} p_i \log_2 p_i \quad \text{bits}.
\]
Shannon proved the **Source Coding Theorem**: the average length $L_{\text{avg}}$ of any uniquely decodable binary code for $X$ satisfies:
\[
  H(X) \;\le\; L_{\text{avg}} \;<\; H(X) + 1.
\]
The name "entropy" was suggested by mathematician John von Neumann, who noted that the mathematical form of $H(X)$ is identical to statistical entropy in thermodynamics.
\end{namebox}

\newpage

% ============================================================
\section{Part 5 \texorpdfstring{---}{--} Solved Exercises}
% ============================================================

\begin{exercise}[Entropy of a Biased Coin]
A biased coin lands heads with probability $p$ and tails with probability $1-p$. Plot the entropy of this source as a function of $p$, and find the value of $p$ that maximizes uncertainty.
\end{exercise}

\textbf{Solution.}
The source $X$ has two outcomes. Its entropy (known as the binary entropy function $H_b(p)$) is:
\[
  H_b(p) \;=\; -p \log_2 p - (1-p) \log_2 (1-p).
\]
To find the maximum, we differentiate $H_b(p)$ with respect to $p$ (using $\frac{d}{dp}\log_2 x = \frac{1}{x \ln 2}$):
\[
  \frac{d}{dp} H_b(p) \;=\; -\log_2 p - \frac{p}{p \ln 2} + \log_2(1-p) + \frac{1-p}{(1-p)\ln 2} \;=\; \log_2 \left(\frac{1-p}{p}\right).
\]
Setting the derivative to 0:
\[
  \log_2 \left(\frac{1-p}{p}\right) \;=\; 0 \implies \frac{1-p}{p} \;=\; 1 \implies 1-p \;=\; p \implies p \;=\; 0.5.
\]
At $p = 0.5$, $H_b(0.5) = -0.5(-1) - 0.5(-1) = 1.0$ bit. The entropy is 0 at $p=0$ and $p=1$, and is symmetric with a maximum at $p=0.5$. This matches intuition: a completely fair coin offers maximum uncertainty (1 bit of information per flip), while a double-headed coin ($p=1$) carries zero uncertainty.

\begin{exercise}[Constructing a Huffman Code]
Construct the optimal prefix-free code for the alphabet $\Sigma = \{A, B, C, D, E\}$ with probabilities:
\[
  p(A) = 0.4, \quad p(B) = 0.2, \quad p(C) = 0.2, \quad p(D) = 0.1, \quad p(E) = 0.1.
\]
Calculate its average length and verify that it satisfies the Source Coding Theorem.
\end{exercise}

\textbf{Solution.}
1. **Combine lowest probabilities:** $D$ (0.1) and $E$ (0.1) are combined into a meta-symbol $DE$ with probability $0.2$. Our set is now $\{A: 0.4, B: 0.2, C: 0.2, DE: 0.2\}$.
2. **Combine next lowest:** Combine $C$ (0.2) and $DE$ (0.2) into $CDE$ with probability $0.4$. Our set is now $\{A: 0.4, B: 0.2, CDE: 0.4\}$.
3. **Combine next lowest:** Combine $B$ (0.2) and $CDE$ (0.4) into $BCDE$ with probability $0.6$. Our set is now $\{A: 0.4, BCDE: 0.6\}$.
4. **Final combination:** Combine $A$ and $BCDE$ into the root.
Now we assign bits from root down:
* $A \mapsto 0$
* $BCDE \mapsto 1$
  * $B \mapsto 10$
  * $CDE \mapsto 11$
    * $C \mapsto 110$
    * $DE \mapsto 111$
      * $D \mapsto 1110$
      * $E \mapsto 1111$

Codeword lengths are: $\ell(A)=1$, $\ell(B)=2$, $\ell(C)=3$, $\ell(D)=4$, $\ell(E)=4$.
Average length:
\[
  L_{\text{avg}} \;=\; 0.4(1) + 0.2(2) + 0.2(3) + 0.1(4) + 0.1(4) \;=\; 0.4 + 0.4 + 0.6 + 0.4 + 0.4 \;=\; 2.2 \text{ bits}.
\]
Now we compute the entropy $H(X)$:
\[
  H(X) \;=\; -0.4\log_2 0.4 - 2(0.2\log_2 0.2) - 2(0.1\log_2 0.1) \;\approx\; 0.528 + 0.929 + 0.664 \;=\; 2.121 \text{ bits}.
\]
Since $2.121 \le 2.2 < 3.121$, the Source Coding Theorem is satisfied.

\bigskip

% ============================================================
\section{Part 6 \texorpdfstring{---}{--} Problems}
% ============================================================

\begin{problem}[Kraft Inequality Proof]
Prove the converse of Kraft's inequality: if a set of positive integers $\ell_1, \ell_2, \ldots, \ell_n$ satisfies $\sum_{i=1}^n 2^{-\ell_i} \le 1$, then there exists a prefix-free binary code with these codeword lengths.
\end{problem}

\begin{problem}[Entropy of a Functions]
Let $X$ be a discrete random variable, and let $g$ be a deterministic function. Prove that:
\[
  H(g(X)) \;\le\; H(X).
\]
Under what conditions does equality hold? Interpret this in terms of processing information.
\end{problem}

\bigskip

% ============================================================
\section{Part 7 \texorpdfstring{---}{--} Deep Problems}
% ============================================================

\begin{deepproblem}[Landauer's Principle and the Physics of Information]
Information is not abstract; it must be stored in physical systems (like voltage levels, electron spins, or molecular configurations).
In 1961, Rolf Landauer proved that erasing a bit of information must dissipate a minimum amount of heat. Let us model this physically.
Consider a single particle in a symmetric double-well potential (representing a 1-bit memory: left well is \textbf{0}, right well is \textbf{1}), in contact with a heat bath at temperature $T$.
\begin{enumerate}[label=(\alph*)]
  \item Initially, the particle is in either well with equal probability (representing an unknown bit). What is the Shannon entropy of this state? What is the thermodynamic entropy $S = k_B \ln \Omega$?
  \item "Reset to Zero" is the process of forcing the particle into the left well (\textbf{0}) regardless of where it started. Show that this process reduces the volume of accessible phase space by a factor of 2.
  \item Since the total entropy of the universe (system + bath) cannot decrease (Second Law of Thermodynamics), show that the entropy of the heat bath must increase by at least $\Delta S_{\text{bath}} \ge k_B \ln 2$.
  \item Prove that this erasure requires an dissipation of heat $Q$ into the bath satisfying:
  \[
    Q \;\ge\; k_B T \ln 2.
  \]
  \item Calculate this heat limit at room temperature ($T = 293$ K) in Joules. Reflect on why modern computers operate at about $10^4$ times this limit, and what physical engineering barriers prevent us from reaching it.
\end{enumerate}
\end{deepproblem}

\bigskip

% ============================================================
\section{Part 8 \texorpdfstring{---}{--} Reflection}
% ============================================================

\begin{reflectionbox}
\textbf{Conceptual Reflection}

Consider the following questions:
\begin{enumerate}
  \item **On physical information:** We often think of software as completely separate from hardware. How does Landauer's Principle change your view of "abstract computation"?
  \item **On structural limits:** In the Kraft inequality, we see that short codes "use up" code space. If a computer's memory addresses must be prefix-free, how does this limit addressable memory size?
\end{enumerate}
\end{reflectionbox}

\vspace{1.4cm}
\begin{center}
{\color{CompBlue}\rule{0.5\linewidth}{0.8pt}}\\[6pt]
{\small\color{DarkGrey}\textit{End of Journey 1, Module 1 --- The Nature of Information}}\\[4pt]
{\small\color{DarkGrey}Continue to: \textbf{Module 1.2 --- Noisy Channels and Error Correction}}\\[6pt]
{\color{CompBlue}\rule{0.5\linewidth}{0.8pt}}
\end{center}

\end{document}
